\chapter{Cross-Cutting Analysis}
\label{ch:cross}
\draftnote{Target length: about 5 pages.}

\section{The Consolidated Comparison}
\label{sec:cross-tables}
% Table 1 — Models and mappings (columns 1–6)
\begin{table}[htbp]
\centering
\small
\caption{Comparison of approaches: models and mappings}
\label{tab:models}
\renewcommand{\arraystretch}{1.3}
\begin{tabularx}{\textwidth}{>{\raggedright\arraybackslash}p{2.0cm}
                             *{5}{>{\raggedright\arraybackslash}X}}
\toprule
\textbf{Approach} & \textbf{Domain \& enforcement point} & \textbf{Abstract model}
  & \textbf{Concrete model} & \textbf{Mapping mechanism} & \textbf{Properties checked} \\
\midrule
HPol \cite{alkhorem2026graph,steiner2018hardening}
  & Web application + DBMS & \emph{TBD} & \emph{TBD} & \emph{TBD} & \emph{TBD} \\
SCM \cite{basile2026scm}
  & Network security controls & \emph{TBD} & \emph{TBD} & \emph{TBD} & \emph{TBD} \\
SMT / ABAC overlay \cite{sissodiya2025k8s}
  & Kubernetes API server & \emph{TBD} & \emph{TBD} & \emph{TBD} & \emph{TBD} \\
\bottomrule
\end{tabularx}
\end{table}

% Table 2 — Evidence and guarantees (key column + columns 7–10)
\begin{table}[htbp]
\centering
\small
\caption{Comparison of approaches: evidence and guarantees}
\label{tab:evidence}
\renewcommand{\arraystretch}{1.3}
\begin{tabularx}{\textwidth}{>{\raggedright\arraybackslash}p{2.0cm}
                             *{4}{>{\raggedright\arraybackslash}X}}
\toprule
\textbf{Approach} & \textbf{Evidence regime} & \textbf{Mechanisation}
  & \textbf{Semantics-preservation claim} & \textbf{Structural limits} \\
\midrule
HPol \cite{alkhorem2026graph,steiner2018hardening}
  & Scenario enumeration; one manual case study
  & Python implementation
  & \emph{TBD (not made?)} & \emph{TBD} \\
SCM \cite{basile2026scm}
  & Expert review, probe packets, attack simulation
  & CLIPS reasoning engine + XML catalogue
  & \emph{TBD (asserted without proof?)} & \emph{TBD} \\
SMT / ABAC overlay \cite{sissodiya2025k8s}
  & Solver counterexample search (sat/unsat)
  & Z3
  & \emph{TBD (quote exact wording for $\tau$)} & \emph{TBD} \\
\bottomrule
\end{tabularx}
\end{table}


\section{Evidence Regimes Compared}
\label{sec:cross-evidence}
\draftnote{Exhaustive scenario enumeration, manual case-study inspection, expert review with
attack simulation, and solver-based counterexample search. For each: the error class it
catches, the error class it structurally cannot catch, and what it costs to run.}

\section{Expressiveness, Decidability and Cost}
\label{sec:cross-tension}
\draftnote{The three-way tension. HPol is decidable and cheap but cannot express context;
SCM is expressive but rests on human judgment; the SMT approach is expressive and
checkable but costs encoding effort and solver time. The NP-hardness of redundancy
elimination, cited in \cite{alkhorem2026graph}, bounds the cheap end.}

\section{The Common Gap}
\label{sec:cross-gap}
\draftnote{Four descents across abstraction levels; four distinct evidence regimes; zero
refinement proofs.}
